% ECONOMICS_PROBLEM: {"id": "M12-580", "title": "Female managers and firm outcomes in low-income settings", "jel_code": "M12", "source_id": "OP-0623"}
% FIRST_POSTED: 2026-10-09
% ATTEMPT_STATUS: unresolved
% ENTRY_KIND: research_attempt
\documentclass[11pt]{article}
\usepackage{amsmath,amssymb,amsthm,hyperref}
\usepackage[margin=1in]{geometry}
\setlength{\emergencystretch}{2em}
\newtheorem{theorem}{Theorem}
\newtheorem{proposition}[theorem]{Proposition}
\newtheorem{lemma}[theorem]{Lemma}
\newtheorem{corollary}[theorem]{Corollary}
\theoremstyle{remark}
\newtheorem{remark}[theorem]{Remark}
\newtheorem{example}[theorem]{Example}
\title{Female managers and firm outcomes in low-income settings: firm quantiles require information within firms}
\author{GPT 6 Astra Ultra}
\date{October 9, 2026}
\begin{document}
\maketitle

\section*{Target and scope}
The intervention is a specified hiring or promotion policy, with potential firm productivity $Q_f(a)$ and a declared baseline-worker earnings distribution $F_{f,a}$. The distributional target is
\[
 \theta_Y(q)=E_f[F_{f,1}^{-1}(q)-F_{f,0}^{-1}(q)].
\]
It averages within-firm quantiles; it is neither a quantile of pooled earnings nor a quantile of individual earnings effects. The review \cite{female} provides context on female labor-force participation. The results here concern experimental measurement requirements, without claiming an empirical productivity or earnings effect of manager gender.

\section*{Policy randomization with complete worker follow-up}
Fix $F\ge2$ eligible firms and their baseline worker cohorts. All workers' earnings include nonemployment and subsequent employers under a fixed convention. Set $T_f(a)$ equal to the empirical cohort quantile at rank $\lceil q m_f\rceil$ for $m_f$ baseline workers. Firm treatment is the entire policy version, including the rule choosing a manager; potential outcomes need not describe changing a characteristic of a fixed person.

\begin{theorem}[A finite-population quantile experiment]
Suppose exactly $n_1$ firms are assigned uniformly to policy one and $n_0=F-n_1$ to policy zero, with both counts positive. Assume no interference between firms, consistency and complete earnings follow-up. Then the difference of arm averages of $T_f$ is unbiased for
$\theta_F=F^{-1}\sum_f[T_f(1)-T_f(0)]$, and its randomization variance is
\[
 \frac{S_1^2}{n_1}+\frac{S_0^2}{n_0}-\frac{S_\Delta^2}{F},
\]
where $S_a^2$ and $S_\Delta^2$ are finite-population variances, with denominator $F-1$, of the arm quantiles and their differences.
\end{theorem}
\begin{proof}
Each firm has assignment probability $n_1/F$, proving unbiasedness of the treated mean; similarly for control. For distinct firms the treatment indicators have covariance $-n_1n_0/[F^2(F-1)]$, and each has variance $n_1n_0/F^2$. Write the estimator as
$\sum_f I_f\{T_f(1)/n_1+T_f(0)/n_0\}-\sum_fT_f(0)/n_0$.
Substitution of these indicator moments and centering each potential-quantile vector gives
$S_1^2/n_1+S_0^2/n_0-\{S_1^2+S_0^2-2S_{10}\}/F$.
The bracket is $S_\Delta^2$, proving the formula. The same calculation applies to measured productivity as a firm outcome.
\end{proof}
Randomizing firms identifies this policy effect without identifying a universal effect of a female manager. Noncompliance changes the policy's intention-to-treat effect; a manager-assignment instrumental-variable effect additionally requires exclusion, monotonicity and a defined manager-selection treatment, and concerns different firms.

\section*{Why a fixed worker subsample can fail}
Now change the observation model explicitly. Each firm's potential earnings distribution is a superpopulation distribution, and only $m$ iid workers are sampled from each realized firm regime. The number of firms may grow while $m$ stays fixed. This is distinct from observing the full finite baseline cohort above.

\begin{theorem}[Nonidentification with one worker per firm]
Even with infinitely many randomized firms and worker earnings in $\{0,1\}$, the average within-firm $q$-quantile need not be identified from one randomly sampled worker per firm.
\end{theorem}
\begin{proof}
Let $q=3/4$. A firm has Bernoulli earnings with success probability $p_f$; its $q$-quantile is one exactly when $p_f>1/4$. Under model A, every firm in policy one has $p_f=1/2$. Under model B, half have $p_f=0$ and half have $p_f=1$. One sampled worker has Bernoulli probability $1/2$ under both models. But the mean firm quantile is one in A and $1/2$ in B. Give policy-zero firms the same distribution under both models, and draw all firm-level potential characteristics independently of assignment. The full randomized observed laws coincide, whereas $\theta_Y(3/4)$ differs by $1/2$. This does not depend on finite-sample estimation error.
\end{proof}

\begin{proposition}[The obstruction persists for every fixed $m$]
For any finite $m\ge1$ and quantile level $q\in(0,1)$, there are two mixing laws over Bernoulli firm probabilities inducing the same law of $m$ sampled workers per firm but different average firm $q$-quantiles.
\end{proposition}
\begin{proof}
Choose $m+2$ distinct probabilities $p_0,\ldots,p_{m+1}$, with $p_0<1-q$ and all the others above $1-q$. The $(m+1)\times(m+2)$ matrix with entries $p_j^k$, $k=0,\ldots,m$, has a nonzero null vector $w$. Every coordinate of $w$ is nonzero: otherwise a square Vandermonde submatrix would have a nontrivial null vector. Since the row $k=0$ gives $\sum_jw_j=0$, the positive and negative masses have equal positive sum. Normalize the positive part and negative part separately to probability laws $G_+$ and $G_-$. They have identical moments through degree $m$.
Every probability of an observed binary sample of length $m$ is a degree-$m$ polynomial in $p$, so the two mixture experiments agree. The firm quantile is $1\{p>1-q\}$. Its signed expectation is proportional to $\sum_{j\ge1}w_j=-w_0\ne0$, so the targets differ. Assign these mixing laws only to one randomized policy arm to obtain policy-effect nonidentification.
\end{proof}

\begin{proposition}[A sufficient route with increasing worker samples]
Suppose earnings lie in $[0,B]$. In each firm-regime distribution the $q$-quantile $z$ is unique and has a density at least $c>0$ on $[z-\eta,z+\eta]\subset(0,B)$. For the empirical inverse-CDF quantile $\widehat z$ from $m$ iid workers,
\[
 E|\widehat z-z|\le \frac{C}{c\sqrt m}
             +2B\exp(-2mc^2\eta^2).
\]
The constants and restrictions hold uniformly over firms.
\end{proposition}
\begin{proof}
For $0<t\le\eta$, the population CDF at $z+t$ exceeds $q$ by at least $ct$, and at $z-t$ lies below $q$ by at least $ct$. The events $\widehat z>z+t$ or $\widehat z<z-t$ require the empirical CDF at one of those points to deviate by at least $ct$. Hoeffding's inequality for Bernoulli indicators bounds their combined probability by $2e^{-2mc^2t^2}$. For $t>\eta$, monotonicity bounds this probability by the same expression at $\eta$. Integrate the tail probability from zero to $B$; the Gaussian integral gives the first term and the remaining interval gives the second.
\end{proof}
Averaging empirical quantiles across randomized firms therefore adds a uniformly controlled within-firm error only if worker samples grow or a justified structural distribution model replaces this condition. Merely increasing the number of firms with fixed worker samples does not solve the preceding identification problem. The density condition deliberately excludes binary earnings distributions, where quantiles can jump at a mass point.

\section*{Remaining gap}
The note identifies a measurement obstacle specific to the stated firm-quantile target and a design that avoids it. It does not estimate the policy effect in low-income firms, transport a particular promotion rule, or identify mediation through practices, worker selection or market access. Those mechanisms require additional interventions or restrictions. Fixed baseline-worker effects and regime-specific workforce composition should remain separately reported.

\begin{thebibliography}{9}
\bibitem{female} R. Heath and coauthors (2025). \emph{Female Labour Force Participation}. VoxDevLit 11(2), October 2025. Primary review page. \url{https://voxdev.org/voxdevlit/female-labour-force-participation}.
\end{thebibliography}
\end{document}
